\documentclass[14pt,a4paper]{article}
\usepackage[utf8]{inputenc}
\usepackage[russian]{babel}
\usepackage{amsmath,amssymb}
\usepackage{geometry}
\geometry{top=2cm,bottom=2cm,left=2.5cm,right=1.5cm}
\usepackage{graphicx}
\usepackage{float}
\usepackage{hyperref}
\hypersetup{
    colorlinks=true,
    linkcolor=blue,
    urlcolor=blue,
    citecolor=blue
}
\usepackage{booktabs}
\usepackage{longtable}
\linespread{1.3}


\begin{document}

\thispagestyle{empty}
\begin{center}
    {\large Санкт-Петербургский государственный университет\par}
    \vspace{0.3cm}
    {\large Прикладная математика и информатика\par}
    \vspace{2.5cm}
    {\large \bfseries ОТЧЕТ\par}
    \vspace{0.5cm}
    {\large по учебной практике 2\par}
    {\large (семестр 2)\par}
    \vspace{0.5cm}
    {\large Сравнение результатов нахождения минимального пути в графе методами Дейкстры и Краскала с применением средств распараллеливания MPI.\par}
\end{center}
\vspace{1.5cm}
\hspace{9cm}{\large Выполнил:\par
    \vspace{0.3cm}
    \hspace{9cm}Олейник Георгий Николаевич,\par
    \hspace{9cm}группа 25.Б22\par
    \vspace{1.7cm}
    \hspace{7cm}Научный руководитель:\par
    \vspace{0.3cm}
    \hspace{7cm}Старший преподаватель мат.-мех. факультета\par
    \hspace{7cm}Мирошниченко Ирина Дмитриевна.\par
    \vspace{0.3cm}
    \hspace{7cm}Кафедра параллельных алгоритмов\par
    \vspace{1.7cm}
    \hspace{6.6cm}Работа выполнена полностью и в срок.\par
    \hspace{6.6cm}Заслуживает оценки A\par
}
\vspace{2cm}
\begin{center}
    {\large Санкт-Петербург\par 2026\par}
\end{center}
\newpage

\begin{center}
    
    {\large \bfseries ОТЗЫВ\par}
    
    \vspace{0.3cm}
    
    на научно-исследовательскую работу\par
    
    \vspace{0.1cm}
    
    \guillemetleft Сравнение результатов нахождения минимального пути в графе методами Дейкстры и Краскала с применением средств распараллеливания MPI\guillemetright\par
    
    \vspace{0.1cm}
    
    студента Олейника Г.Н. (группа 25.Б22)\par
    
    \vspace{2cm}

\end{center}


\noindent \hspace{2.4cm} Представленная тема работы является актуальной. При выполнении работы

\noindent \hspace{2.4cm} студент проявил самостоятельность, получил хорошие результаты. Работа

\noindent \hspace{2.4cm} выполнена полностью и в срок, все поставленные задачи выполнены.

\noindent \hspace{2.4cm} Необходимые компетенции согласно программе дисциплины и

\noindent \hspace{2.4cm} образовательной программы сформированы. Оценка отлично (А).

\vspace{1.7cm}

\hspace{-1.55cm}Старший преподаватель мат.-мех. факультета, \hspace{6.5cm} Мирошниченко И.Д.\par
\hspace{-1.55cm}03.06.2026

\newpage
\tableofcontents
\newpage

\section{Введение}
Современные вычислительные системы всё чаще строятся по принципу многоядерных и многопроцессорных архитектур. Это связано с тем, что дальнейшее увеличение тактовой частоты одного ядра стало наталкиваться на физические ограничения (тепловыделение, утечки тока). Поэтому производительность теперь наращивают за счёт увеличения числа ядер, объединяемых в одном процессоре, а также за счёт использования многопроцессорных систем и кластеров.

Однако просто наличие нескольких ядер не даёт выигрыша — необходимо адаптировать классические алгоритмы к параллельному исполнению, уметь разбивать задачу на подзадачи, распределять их между вычислительными элементами и затем собирать результат. Для систем с распределённой памятью (например, кластеров) одним из основных инструментов является стандарт MPI (Message Passing Interface). MPI предоставляет программисту набор функций для обмена сообщениями между процессами, коллективных операций и синхронизации.

В данной работе рассматриваются два фундаментальных алгоритма на графах — \textbf{алгоритм Дейкстры} (поиск кратчайших путей от источника до всех вершин) и \textbf{алгоритм Краскала} (построение минимального остовного дерева). Эти алгоритмы широко применяются в задачах маршрутизации (построение кратчайших маршрутов в транспортных сетях), кластеризации (выделение связных компонент), проектировании телекоммуникационных сетей, а также в биоинформатике (построение эволюционных деревьев).

\section{Постановка задачи}
Необходимо решить следующие задачи:
\begin{enumerate}
    \item Разработать последовательные программы на C++ для алгоритма Дейкстры и алгоритма Краскала.
    \item Разработать параллельные программы для этих же алгоритмов с использованием технологии MPI.
    \item Сгенерировать тестовые неориентированные графы с плотностью 20\% и весами от 1 до 100 для размеров $n = 500, 1000, 5000, 10000$ вершин.
    \item Провести вычислительные эксперименты: измерить время выполнения последовательных версий и параллельных версий при количестве процессов $p = 1, 2, 4, 8$.
    \item Вычислить ускорение $S_p = T_{\text{seq}} / T_p$ и эффективность $E_p = S_p / p$ для каждого размера графа и числа процессов.
    \item Сравнить поведение двух алгоритмов, выявить причины отсутствия или наличия ускорения.
    \item Сформулировать выводы о применимости распараллеливания для графов среднего размера.
\end{enumerate}

\section{Основные определения}
\subsection{Графы}
Граф $G = (V,E)$ состоит из множества вершин $V$ и множества рёбер $E$. Если граф взвешенный, то каждому ребру $(u,v)$ приписан вес $w(u,v)$, обычно положительное число. Кратчайшим путём из вершины $s$ в вершину $t$ называется путь с минимальной суммой весов среди всех возможных путей между этими вершинами. Минимальное остовное дерево (MST) — это связный подграф, содержащий все вершины исходного графа, являющийся деревом (не содержит циклов) и имеющий минимальную суммарную массу.

\subsection{Технология MPI}
MPI (Message Passing Interface) — это стандарт и библиотека, предоставляющая функции для организации параллельных вычислений в системах с распределённой памятью. Основные понятия:
\begin{itemize}
    \item \textbf{Процесс} — независимо выполняемый экземпляр программы, имеющий уникальный номер (ранг) в рамках коммуникатора.
    \item \textbf{Коммуникатор} — группа процессов, которые могут обмениваться сообщениями; стандартный коммуникатор, включающий все запущенные процессы, называется \texttt{MPI\_COMM\_WORLD}.
    \item \textbf{Коллективные операции} — взаимодействие всех процессов коммуникатора (рассылка \texttt{MPI\_Bcast}, сбор данных \texttt{MPI\_Gather}, глобальные вычисления \texttt{MPI\_Allreduce}).
    \item \textbf{Операции точка-точка} — обмен данными между двумя процессами (\texttt{MPI\_Send}, \texttt{MPI\_Recv}).
    \item \textbf{Барьерная синхронизация} (\texttt{MPI\_Barrier}) — останавливает выполнение всех процессов до тех пор, пока каждый из них не вызовет эту функцию.
\end{itemize}
Правильное использование MPI позволяет добиться высокой производительности на многопроцессорных системах, однако требует тщательного учёта баланса между вычислениями и коммуникациями.

\section{Алгоритмы и их последовательные версии}
\subsection{Алгоритм Дейкстры}
Алгоритм Дейкстры (предложен Э. Дейкстрой в 1959 году) решает задачу о кратчайших путях из одной вершины во взвешенном графе с неотрицательными весами. Классическая последовательная реализация использует очередь с приоритетом (например, бинарную кучу). Сложность составляет $O((V+E)\log V)$. Для плотных графов ($E \approx V^2$) это $O(V^2 \log V)$.

Псевдокод:
\begin{verbatim}
dist[s] = 0; остальным dist = INF
pq.push((0, s))
while pq не пуста:
    (d, u) = pq.pop()
    if d > dist[u]: continue
    for (v, w) in g[u]:
        if dist[v] > d + w:
            dist[v] = d + w
            pq.push((dist[v], v))
\end{verbatim}
Алгоритм работает следующим образом: из очереди извлекается вершина с наименьшим текущим расстоянием, после чего проверяются все её соседи; если удаётся улучшить расстояние до соседа через текущую вершину, то расстояние обновляется и вершина помещается в очередь.

\subsection{Алгоритм Краскала}
Алгоритм Краскала (Дж. Краскал, 1956) строит минимальное остовное дерево. Сначала все рёбра сортируются по весу, затем добавляются в MST, если они не образуют цикла (проверка через систему непересекающихся множеств DSU). Сложность $O(E \log E)$ за счёт сортировки, а затем почти линейный проход. Для плотных графов это также $O(V^2 \log V)$.

Псевдокод:
\begin{verbatim}
отсортировать рёбра по весу
инициализировать DSU (n)
для каждого ребра (u,v,w) в порядке возрастания:
    if find(u) != find(v):
        добавить ребро в MST
        union(u,v)
        if |MST| == n-1: break
\end{verbatim}
DSU позволяет быстро проверять, принадлежат ли две вершины одной компоненте, и объединять компоненты. Алгоритм завершается, когда добавлено $V-1$ ребро.

\section{Параллельные реализации на MPI}
\subsection{Параллельная версия Дейкстры}
В параллельной версии вершины распределяются между процессами блочно. Каждый процесс хранит локальные массивы расстояний для своих вершин и флаги посещения. Для локального поиска минимума используется своя куча. Основной цикл (на каждой итерации):
\begin{enumerate}
    \item Каждый процесс извлекает вершину с минимальным расстоянием из своей кучи (локальный минимум).
    \item С помощью \texttt{MPI\_Gather} все локальные минимумы (значение и глобальный индекс вершины) собираются на процессе 0.
    \item Процесс 0 выбирает глобальный минимум и рассылает его всем остальным процессам через \texttt{MPI\_Bcast}.
    \item Каждый процесс, получив глобальную вершину и её расстояние, выполняет релаксацию рёбер из этой вершины, но только для тех вершин, которые ему принадлежат.
    \item Процесс, которому принадлежит выбранная вершина, помечает её как посещённую и удаляет из своей кучи.
\end{enumerate}
Недостаток этого подхода — необходимость синхронизации на каждой из $V$ итераций. Даже если сама релаксация выполняется параллельно, накладные расходы на коллективные операции растут с числом процессов и могут перекрыть выигрыш.

\subsection{Параллельная версия Краскала}
Параллельная версия Краскала основана на распределении рёбер между процессами:
\begin{enumerate}
    \item Процесс 0 генерирует все рёбра графа и сортирует их глобально (используя быструю сортировку).
    \item Отсортированный список рёбер разбивается на непрерывные блоки, и каждый процесс получает свой блок (через \texttt{MPI\_Send} / \texttt{MPI\_Recv} или коллективные операции).
    \item Каждый процесс независимо строит локальное MST, применяя DSU только к своей порции рёбер. Это даёт набор рёбер, которые являются остовом для подмножества вершин (локальный лес).
    \item Локальные списки рёбер (размер каждого не превышает $V-1$) собираются на процессе 0.
    \item Процесс 0 выполняет финальное построение глобального MST: сортирует собранные рёбра (они уже почти упорядочены) и применяет DSU, чтобы объединить компоненты.
\end{enumerate}
Основные узкие места: глобальная сортировка на одном процессе и пересылка рёбер. Однако локальное построение MST можно выполнять параллельно, что даёт некоторое ускорение на больших графах.

\section{Реализация программ}
Полные исходные коды (последовательные и параллельные версии для обоих алгоритмов) размещены в открытом репозитории GitHub:

\url{https://github.com/bolgarka38/research}


Все программы написаны на C++ с использованием MS-MPI. Для компиляции используется команда \texttt{mpicxx /EHsc имя\_файла.cpp -o имя\_файла.exe}. Для запуска — \texttt{mpiexec -n N имя\_файла.exe размер\_графа}. Для последовательных версий компиляция и запуск выполняются без MPI.

\section{Экспериментальная установка}
Программы компилировались с оптимизацией /O2 (MSVC). Вычислительная система: персональный компьютер с 4 физическими ядрами (8 логических процессоров), операционная система Windows 10. Время измерялось с помощью \texttt{MPI\_Wtime} (параллельные версии) и std::chrono (последовательные). Каждый эксперимент повторялся 5 раз, в таблицах приведены средние значения. Тестовые графы: неориентированные, плотность 20\%, веса рёбер — случайные целые числа от 1 до 100, размеры $n = 500, 1000, 5000, 10000$. Графы генерируются с фиксированным начальным зерном (seed), что обеспечивает воспроизводимость результатов.

\section{Результаты экспериментов}
\subsection{Результаты для алгоритма Дейкстры}
В таблице \ref{tab:dijkstra} приведены время последовательной версии ($T_{\text{seq}}$), время MPI-версии ($T_p$) для $p = 1,2,4,8$, ускорение $S_p = T_{\text{seq}} / T_p$ и эффективность $E_p = S_p / p$. Обратите внимание, что даже при $p=1$ MPI-версия заметно медленнее последовательной из-за накладных расходов на инициализацию и вызовы функций.

\begin{table}[H]
\centering
\caption{Результаты для алгоритма Дейкстры (heap-версия)}
\label{tab:dijkstra}
\begin{tabular}{|c|c|c|c|c|c|}
\hline
$n$ & $T_{\text{seq}}$, с & $p$ & $T_p$, с & $S_p$ & $E_p$ \\
\hline
500 & 0.0007805 & 1 & 0.0015165 & 0.515 & 0.515 \\
500 & 0.0007805 & 2 & 0.0022385 & 0.349 & 0.174 \\
500 & 0.0007805 & 4 & 0.0039459 & 0.198 & 0.049 \\
500 & 0.0007805 & 8 & 0.0048767 & 0.160 & 0.020 \\
\hline
1000 & 0.0021731 & 1 & 0.0047562 & 0.457 & 0.457 \\
1000 & 0.0021731 & 2 & 0.0068324 & 0.318 & 0.159 \\
1000 & 0.0021731 & 4 & 0.0125590 & 0.173 & 0.043 \\
1000 & 0.0021731 & 8 & 0.0209469 & 0.104 & 0.013 \\
\hline
5000 & 0.0348004 & 1 & 0.0855380 & 0.407 & 0.407 \\
5000 & 0.0348004 & 2 & 0.1108900 & 0.314 & 0.157 \\
5000 & 0.0348004 & 4 & 0.1126010 & 0.309 & 0.077 \\
5000 & 0.0348004 & 8 & 0.1671630 & 0.208 & 0.026 \\
\hline
10000 & 0.1327200 & 1 & 0.3410300 & 0.389 & 0.389 \\
10000 & 0.1327200 & 2 & 0.4038480 & 0.329 & 0.164 \\
10000 & 0.1327200 & 4 & 0.5289730 & 0.251 & 0.063 \\
10000 & 0.1327200 & 8 & 0.7165500 & 0.185 & 0.023 \\
\hline
\end{tabular}
\end{table}

Из таблицы видно, что время выполнения MPI-версии Дейкстры растёт с увеличением числа процессов. Ускорение $S_p$ не превышает 0.5, а эффективность $E_p$ при $p=8$ падает до 0.02–0.023, что говорит о крайне низкой масштабируемости.

\subsection{Результаты для алгоритма Краскала}
Таблица \ref{tab:kruskal} содержит аналогичные данные. Для $n=10000$ последовательное время $T_{\text{seq}}=0.87641$ с (измерено отдельно). Обратите внимание, что абсолютное ускорение $S_p$ остаётся ниже 0.34, но при переходе от 1 к 8 процессам время $T_p$ уменьшается.

\begin{table}[H]
\centering
\caption{Результаты для алгоритма Краскала}
\label{tab:kruskal}
\begin{tabular}{|c|c|c|c|c|c|}
\hline
$n$ & $T_{\text{seq}}$, с & $p$ & $T_p$, с & $S_p$ & $E_p$ \\
\hline
500 & 0.0023026 & 1 & 0.0074298 & 0.310 & 0.310 \\
500 & 0.0023026 & 2 & 0.0076061 & 0.303 & 0.151 \\
500 & 0.0023026 & 4 & 0.0079199 & 0.291 & 0.073 \\
500 & 0.0023026 & 8 & 0.0082092 & 0.280 & 0.035 \\
\hline
1000 & 0.0089040 & 1 & 0.0304275 & 0.293 & 0.293 \\
1000 & 0.0089040 & 2 & 0.0351443 & 0.253 & 0.127 \\
1000 & 0.0089040 & 4 & 0.0305284 & 0.292 & 0.073 \\
1000 & 0.0089040 & 8 & 0.0305684 & 0.291 & 0.036 \\
\hline
5000 & 0.2194330 & 1 & 0.7227170 & 0.304 & 0.304 \\
5000 & 0.2194330 & 2 & 0.6777040 & 0.324 & 0.162 \\
5000 & 0.2194330 & 4 & 0.6844660 & 0.321 & 0.080 \\
5000 & 0.2194330 & 8 & 0.6686740 & 0.328 & 0.041 \\
\hline
10000 & 0.8764100 & 1 & 2.8095900 & 0.312 & 0.312 \\
10000 & 0.8764100 & 2 & 2.6957100 & 0.325 & 0.163 \\
10000 & 0.8764100 & 4 & 2.6406300 & 0.332 & 0.083 \\
10000 & 0.8764100 & 8 & 2.6342500 & 0.333 & 0.042 \\
\hline
\end{tabular}
\end{table}

\section{Анализ результатов}
\subsection{Анализ для алгоритма Дейкстры}
Для алгоритма Дейкстры наблюдается монотонный рост времени выполнения с увеличением числа процессов. Это объясняется тем, что на каждой из $n$ итераций выполняются коллективные операции \texttt{MPI\_Gather} и \texttt{MPI\_Bcast}. Даже если каждая операция занимает всего несколько микросекунд, при $n=10000$ суммарные накладные расходы достигают десятков миллисекунд, что сопоставимо с общим временем счёта. Более того, с ростом числа процессов $p$ объём данных, пересылаемых при \texttt{MPI\_Gather} (локальные минимумы), увеличивается, а также растёт задержка на синхронизацию. В результате параллельная версия не только не даёт ускорения, но и работает значительно медленнее последовательной. Даже при $p=1$ MPI-версия медленнее из-за накладных расходов на вызовы функций инициализации и коллективных операций (которые в реализации Microsoft могут быть не оптимизированы для одного процесса). Эффективность $E_p$ падает до 0.023 при $p=8$, что свидетельствует о крайне низкой масштабируемости.

\subsection{Анализ для алгоритма Краскала}
Алгоритм Краскала демонстрирует более интересное поведение. Хотя абсолютное ускорение относительно последовательной версии $S_p$ остаётся ниже 0.34 (то есть MPI-версия всё равно медленнее), относительное ускорение самой MPI-версии при увеличении числа процессов положительно. Для $n=10000$: $T_1=2.810$ с, $T_8=2.634$ с. Относительное ускорение составляет $2.810 / 2.634 \approx 1.067$ (6.6\%). Для $n=5000$: $0.7227 / 0.6687 \approx 1.081$ (8.1\%). Это означает, что распределённая обработка рёбер (локальное построение MST) даёт некоторый выигрыш, который, однако, не перекрывает накладные расходы на глобальную сортировку и финальное слияние. При $n=1000$ и $n=500$ улучшения практически нет, так как доля коммуникаций и последовательной сортировки становится доминирующей.

\subsection{Закон Амдала}
Закон Амдала гласит, что максимальное ускорение $S$ ограничено долей последовательного кода $\sigma$: $S \le 1 / (\sigma + (1-\sigma)/p)$. В параллельной версии Краскала последовательной частью являются глобальная сортировка всех рёбер (выполняется на процессе 0) и финальное слияние. Если принять $\sigma = 0.8$ (80\% времени тратится на сортировку и слияние), то максимальное ускорение на 8 процессах составит $1 / (0.8 + 0.2/8) = 1.19$. Экспериментальные значения (1.067–1.08) находятся в этом диапазоне, что подтверждает ограничения, накладываемые последовательной частью. Для Дейкстры же доля последовательного кода ещё выше, так как каждая итерация требует глобальной синхронизации, что делает параллельную версию практически бесполезной.

\section{Выводы}
На основе проведённых экспериментов можно сделать следующие выводы:
\begin{itemize}
    \item Параллельная реализация Дейкстры с синхронизацией на каждой итерации неэффективна для графов до 10000 вершин. Время выполнения растёт с числом процессов, ускорение $S_p$ не превышает 0.5, эффективность $E_p$ падает до 0.02–0.03. Основная причина — высокие накладные расходы на коллективные операции и малая вычислительная нагрузка на процесс.
    \item Параллельная реализация Краскала даёт небольшое относительное ускорение (6–8\%) при переходе от 1 к 8 процессам на графах от 5000 вершин. Это объясняется распараллеливанием обработки рёбер после глобальной сортировки. Однако абсолютное ускорение относительно последовательной версии остаётся ниже 0.34, то есть MPI-версия всё равно медленнее.
    \item Для получения значительного ускорения необходимы более совершенные параллельные алгоритмы (например, $\Delta$-stepping для Дейкстры, параллельные алгоритмы MST типа алгоритма Борувки) или использование существенно больших графов (порядка сотен тысяч вершин), где вычислительная нагрузка начинает доминировать над коммуникациями.
    \item Последовательные версии обоих алгоритмов остаются предпочтительными для задач среднего размера (до 10000 вершин). При этом, если всё же требуется распараллеливание, для Краскала можно рекомендовать использование 2–4 процессов, так как дальнейшее увеличение даёт лишь незначительный выигрыш.
\end{itemize}

\section{Заключение}
В работе были реализованы последовательные и параллельные (MPI) версии алгоритмов Дейкстры и Краскала. Проведённые эксперименты показали, что прямое («лобовое») распараллеливание не даёт выигрыша для графов среднего размера из-за доминирования коммуникационных накладных расходов. Тем не менее, для Краскала наблюдался слабый положительный эффект, что подтверждает принципиальную возможность распараллеливания при условии большой вычислительной нагрузки. Полученные результаты могут быть использованы при выборе стратегии параллелизации для реальных графовых задач, а также при проектировании более эффективных параллельных алгоритмов. Полные исходные коды доступны в репозитории: \url{https://github.com/bolgarka38/research}.

\section*{Список литературы}
\begin{enumerate}
    \item Антонов А.С. Параллельное программирование с использованием технологии MPI. – М.: Изд-во МГУ, 2004. – 72 с.
    \item Воеводин В.В., Воеводин Вл.В. Параллельные вычисления. – СПб.: БХВ-Петербург, 2002. – 608 с.
    \item Кормен Т. и др. Алгоритмы: построение и анализ. – М.: Вильямс, 2013. – 1328 с.
    \item MPI: A Message-Passing Interface Standard (Version 1.1). – \url{https://www.mpi-forum.org/} – 548 с.
\end{enumerate}
\end{document}